\section{Native causal forests: preserve the target before accelerating it}
\label{exec21:cxx-scope}
A servicing team considering different offers for different risk groups needs more than an accurate forecast of repayment. It needs a contrast between actions, a declared population over which to average that contrast, and uncertainty that corresponds to the decision. A C++ implementation can reduce computation while changing any of these meanings at a language boundary. The useful comparison therefore starts with the statistical object and follows it through wrappers, native learning, stored draws and portfolio aggregation. Four inspected lineages provide complementary designs: generalized random forests, \texttt{bcf}, XBCF and flexBCF. Their common use of trees does not make their targets, priors or intervals interchangeable.\citep{exec21_cxx_grf,exec21_cxx_bcf,exec21_cxx_xbcf,exec21_cxx_flexbcf}

The comparison binds grf 2.6.1 to its 30 April 2026 commit, bcf 2.0.2 to 26 September 2023, XBCF's declared 0.1.2 source to 25 June 2024, and flexBCF 1.0.3 to 1 December 2025. These are inspected versions preceding the literature cutoff, not assertions about the newest release. The native grf boundary analysis developed elsewhere in this monograph remains applicable: the initial proposal favors a qualified R process around its native core, rather than a new direct embedding or port. This comparison extends the alternatives and semantic contracts; it supplies no runtime qualification.\citep{exec21_cxx_grf,exec21_cxx_bcf,exec21_cxx_xbcf,exec21_cxx_flexbcf}

\subsection{A conditional effect and a dose slope answer different questions}
\label{exec21:cxx-target}
For a binary action $A$, let $Y(a)$ be the potential outcome under action $a$ and $X$ permitted baseline information. Under consistency, no interference, conditional ignorability and overlap, the conditional average effect is
\begin{equation}
 \tau(x)=E\{Y(1)-Y(0)\mid X=x\}.
 \label{exec21:cxx-cate}
\end{equation}
The inspected grf wrapper describes this binary target, estimates or accepts aligned outcome and treatment nuisance predictions, and residualizes before forest fitting. Honesty separates split selection from leaf estimation. The native trainer checks whether the sample and honesty fractions leave observations in both parts; its tree trainer splits sampled clusters and repopulates estimation leaves. Computational separation does not establish ignorability or validate a treatment clock.\citep{exec21_cxx_grf}

For continuous treatment the wrapper documents a conditional covariance-to-variance slope,
\begin{equation}
 s(x)=\frac{\operatorname{Cov}(Y,A\mid X=x)}{\operatorname{Var}(A\mid X=x)}.
 \label{exec21:cxx-slope}
\end{equation}
It cannot be substituted automatically for $E\{Y(a_1)-Y(a_0)\mid X=x\}$. In an authored example, randomize $A$ uniformly over $\{-1,0,1\}$ and set $Y(a)=a^2$. The covariance numerator is zero by symmetry, while the action contrast from zero to one is one. Neither confounding nor a computational failure is needed for the discrepancy: the requested objects differ. A bank seeking the response to a specific change in fee or service intensity must request that dose pair or justify a response model under which the slope represents it. This is an exact semantic witness, not a run of grf.\citep{exec21_cxx_grf,exec21_cxx_reference}

The native prediction strategy solves a weighted local moment ratio. Its instrumental implementation divides an outcome--instrument covariance expression by a treatment--instrument expression. A near-zero denominator is a scientific and numerical warning. The inspected variance code groups trees, skips groups with empty leaves and corrects between- and within-group quantities. Its objective-Bayes correction avoids negative variance estimates; that label does not make its output draws from a Bayesian causal forest. An adapter must retain effective group counts, denominator diagnostics and unsupported variance status. Full upstream guard behavior requires a qualified native trial.\citep{exec21_cxx_grf}

\subsection{Why regularizing prediction can distort a treatment contrast}
\label{exec21:cxx-prior}
The genuine inspected BCF original is the 33-page arXiv version 4 of 13 November 2019 by Hahn, Murray and Carvalho; it is not asserted to be byte-identical to the later journal article. Under its stated assumptions the conditional mean is modeled as
\begin{equation}
 E(Y\mid X=x,A=a)=\mu\{x,\widehat\pi(x)\}+a\tau(x),
 \label{exec21:cxx-bcf-model}
\end{equation}
with separate tree priors for the prognostic and treatment functions. The fitted propensity $\widehat\pi$ informs the response-surface prior. This addresses regularization-induced confounding: shrinkage of the prognostic function can leave treatment to explain outcome variation associated with selection. A propensity feature does not imply an orthogonal estimating equation or double robustness. The paper explicitly disclaims double robustness for its natural Bayesian estimates and discusses sensitivity to alternative propensity estimates.\citep{exec21_cxx_hahn}

A project derivation isolates the mechanism. Let $X$ be equally likely to be $-1$ or $1$, with treatment probabilities $1/4$ and $3/4$ respectively. Set $Y(a)=2X$ for both actions, so the true causal effect is zero and overlap holds. Write $Z=A-1/2$. Minimize
\begin{equation}
 E\{(Y-\beta X-tZ)^2\}+\lambda\beta^2,
 \qquad \lambda=3/4.
 \label{exec21:cxx-shrink-objective}
\end{equation}
Only the prognostic coefficient is penalized. Here $E(X^2)=1$, $E(Z^2)=1/4$ and $E(XZ)=1/4$. The treatment normal equation gives $t=2-\beta$; the prognostic equation then gives
\begin{equation}
 \beta=\frac{2(3/4)}{3/4+\lambda}=1,
 \qquad t=1.
 \label{exec21:cxx-shrink-result}
\end{equation}
The apparent treatment effect is produced by the stipulated shrinkage. The exact reference verifies both normal equations. This motivates studying prior design, treatment selection and nuisance learning jointly. It is not a theorem that every regularized learner fails or that BCF necessarily repairs every selection mechanism.\citep{exec21_cxx_reference,exec21_cxx_hahn}

The bcf wrapper exposes separate prognostic and moderation matrices, supplied propensities and their placement in either or both ensembles. Prior-scale choices and outcome standardization belong to the fitted model. A port that copies trees but drops a scale prior changes the procedure. flexBCF explicitly distinguishes its default model from bcf in scale priors, treatment-recoding behavior and relative regularization of the two ensembles. A claim of speed or flexibility therefore cannot establish identical inference. A prototype comparison must freeze each configuration and compare a common target.\citep{exec21_cxx_bcf,exec21_cxx_flexbcf}

\subsection{Three meanings of weight must remain separate}
\label{exec21:cxx-weights}
The bcf argument \texttt{w} represents likelihood precision: the conditional variance is $\sigma^2/w_i$. The native source uses these weights in tree updates and residual sums. They do not define the population for a reported effect. For a declared cohort $\mathcal I$ and target weights $q_i$, a draw-wise average is
\begin{equation}
 T^{(s)}=\frac{\sum_{i\in\mathcal I}q_i\tau_i^{(s)}}{\sum_{i\in\mathcal I}q_i}.
 \label{exec21:cxx-aggregate}
\end{equation}
Likelihood precision, target weights and cohort selection require different fields. The target can concern eligible customers, treated customers or a transported population; its meaning cannot be recovered from an unlabeled vector.\citep{exec21_cxx_bcf,exec21_cxx_flexbcf}

For two equally represented groups with effects $-2$ and $6$, the population effect is $2$. Reusing precisions $1$ and $9$ as target weights gives $26/5=5.2$. If these effects are stipulated dollar gains and the offer costs USD 3 per customer, the population margin is USD -1 while the reweighted margin is USD 2.20. Better observation precision does not authorize replacing the eligible customer mix.\citep{exec21_cxx_reference}

In grf, equalizing cluster weights changes training samples and average-effect weights. The documented training rule draws equal numbers of units from sampled clusters, using the smallest cluster size. This option is incompatible with supplying the ordinary sample-weight vector in that interface. Balancing branches instead of customers changes both the target and parts of learning; it cannot be implemented by relabeling a final average while leaving an incompatible training run untouched.\citep{exec21_cxx_grf}

\subsection{Joint draws, scale restoration and persistence}
\label{exec21:cxx-draws}
A Bayesian portfolio interval averages customer effects within a common posterior draw before summarizing portfolio draws. Averaging individual interval endpoints generally differs. Consider two equally weighted customers with joint draws $(-2,6)$ and $(6,-2)$. Each customer's marginal support is $[-2,6]$, but every portfolio draw is exactly $2$. Averaging marginal endpoints discards the perfect negative dependence. At the assumed USD 3 cost every portfolio margin is USD -1. Two artificial draws illustrate algebra; they establish neither convergence nor calibration.\citep{exec21_cxx_reference}

The inspected flexBCF native SATT routine evaluates each stored ensemble on supplied rows, averages within the draw and then calculates quantiles. Its weighted companion divides by the supplied weight sum. The routine name does not select treated rows: cohort selection must already be reflected in its inputs. These selected functions do not settle how all public callers validate empty cohorts, nonfinite weights, zero sums or category levels. The proposed adapter must require explicit cohort and weight contracts; complete caller tracing remains open.\citep{exec21_cxx_flexbcf}

Scale parameters must remain aligned by draw. In inspected bcf prediction an ensemble output $g_i^{(s)}$ contributes
\begin{equation}
 \tau_i^{(s)}=s_Y\{b_1^{(s)}-b_0^{(s)}\}g_i^{(s)},
 \label{exec21:cxx-restore}
\end{equation}
where $s_Y$ is the outcome scale. A treatment contrast receives no added outcome mean. XBCF also restores a treatment draw using outcome scale and its treatment-scale difference. If raw draws are $1,2$, $s_Y=10$, and scale differences are $1/2,2$, restored draws are $5,40$ with mean $22.5$. Omitting the contrast factor gives mean $15$. At an assumed USD 20 action cost these imply USD 2.50 and USD -5 margins. Multiplying mean raw output by mean scale also loses their posterior dependence.\citep{exec21_cxx_bcf,exec21_cxx_xbcf,exec21_cxx_reference}

bcf documents draws by units, whereas selected XBCF outputs use units by iterations. Transport needs unit, chain and retained-draw identifiers rather than a guessed transpose. bcf's wrapper serializes tree samples for later prediction; suppressing file output prevents that path. Its selected prediction wrapper requires both treatment labels, a restriction to reconcile for a one-cohort request. The chain executor sets one core while retaining within-chain thread controls. These are persistence and execution questions for an adapter, not measured speed or scalability.\citep{exec21_cxx_bcf,exec21_cxx_xbcf}

\subsection{Static hazards define focused prototype tests}
\label{exec21:cxx-hazards}
In pinned XBCF's combined R predictor, the default propensity branch refers to \texttt{z}, which is not a formal argument. The function assigns predictions to \texttt{obj1} and \texttt{obj2}, then reads \texttt{obj} before its local assignment at the end. The separate treatment-draw predictor uses its own local object. These are particular source paths requiring isolated reconciliation, not failures observed by executing the package. An enclosing environment, later release or different entrypoint must not silently replace the inspected route.\citep{exec21_cxx_xbcf}

In the selected Python wrapper, the zero-outcome-variance branch constructs \texttt{ValueError} without raising it, and adjusted fitted draws are assigned inside the standardization branch. Prediction has separate restoration logic. Tests should cover constant outcomes, both standardization settings, fit-versus-predict parity and return formats. The native loop alternates prognostic and treatment updates and updates scales; stored output must be checked against wrapper order. Source resemblance does not prove equivalent R and Python interfaces.\citep{exec21_cxx_xbcf}

The selected flexBCF trainer derives the control count from prognostic and treatment design sizes and offsets treatment updates accordingly. Its wrapper transposes inputs and passes the treatment indicator. Training order and treatment-specific rows therefore require tracing through public callers and documentation. Assuming the same row layout across all four libraries would hide a consequential boundary error.\citep{exec21_cxx_flexbcf}

\subsection{Historical performance supports an experiment}
\label{exec21:cxx-evidence}
The BCF paper offers evidence for its design and relevant limitations. Its targeted-selection example improves bias and coverage over ordinary BART under that simulation; broader experiments do not give uniformly nominal coverage. In the retrospective ACIC 2016 comparison, Table 4 reports nominal 95\% interval coverage of 82\% for BCF, 88\% for propensity-score BART, 81\% for BART and 58\% for the compared causal random forest. BCF's mean interval length is 0.026 versus 0.055 for that forest. Shorter intervals and lower estimation error coexist with undercoverage. The paper discusses this limitation and supplies no general frequentist coverage theorem for its credible intervals.\citep{exec21_cxx_hahn}

These are historical reported results, not a project reproduction or ranking of today's implementations. They motivate public experiments reporting effect error, coverage, width, failures and decision consequences together. Known counterfactual generators can vary overlap, targeted selection, effect sparsity, response nonlinearity, propensity misspecification and dependence. Public benchmarks provide complementary designs whose versions, targets and access conditions must be fixed. Failed fits and invalid requests stay in the denominator. Posterior uncertainty, Monte Carlo error and repeated-sampling coverage must remain separate: more draws do not repair a wrong causal target or unmeasured confounding.\citep{exec21_cxx_hahn,exec21_cxx_reference}

\subsection{Integration choices and the remaining census}
\label{exec21:cxx-integration}
For the proposal, grf remains a candidate through its analyzed R/native boundary. bcf is a candidate for a process-based Bayesian adapter with explicit priors and persisted trees. XBCF warrants a performance comparison after selected wrapper paths are reconciled. flexBCF is a distinct alternative needing prior, coding, training-order and cohort contracts. Porting would transfer responsibility for random numbers, tree updates, priors, draw alignment and persistence. The initial approach should preserve each native model behind a typed adapter and qualify a narrow target before considering a port. These are reasoned proposal choices, not completed reuse admissions.\citep{exec21_cxx_grf,exec21_cxx_bcf,exec21_cxx_xbcf,exec21_cxx_flexbcf}

Inspected declarations are GPL-family terms for grf, bcf and flexBCF and Apache 2.0 for XBCF. They identify a rights review rather than completing one: dependencies, notices, distribution and embedding remain to be examined. A separate process does not automatically settle licensing obligations. Bounded repository discovery also found SparseBCF, time-series, multivariate, survival and other candidates with CRAN mirrors. Names are leads; lineage, cutoff availability, originals and capabilities remain uninspected. Mirrors do not count as independent implementations.\citep{exec21_cxx_grf,exec21_cxx_bcf,exec21_cxx_xbcf,exec21_cxx_flexbcf}

Library, CLI and MCP requests should carry action contrast or slope, cohort, clock, source version, nuisance provenance, weight roles, uncertainty type and refusal reasons. An agent can help formulate requests and inspect diagnostics, but cannot relabel posterior quantiles as calibrated confidence limits. These exact witnesses justify inexpensive public prototype tests before costly bank trials. Broader coverage, method originals, rights, native parity and calibrated decision performance remain prerequisites for a stronger adoption claim.
